% College Lectures 10-11: primes, primality tests, factorization, CRT, fast exponentiation, discrete logs.
% Every number in this file was recomputed by script (scratchpad/primes/verify*.py).
% Fermat/Euler second versions, inverses by Fermat/Euler, order and primitive roots now live in 11_fermat_euler.tex (sec:fermat).
\section{Primes, Primality Testing, Factorization, CRT and Discrete Logs}\label{sec:primes}

Asymmetric-key cryptography is built on a handful of facts about primes, and this chapter collects them. RSA needs two large random primes, so we need a fast way to decide whether a 1024-bit number is prime; this is \emph{primality testing}. RSA is secure only because nobody can split $n = pq$ back into $p$ and $q$; this is the hardness of \emph{factorization}. RSA decryption is made about four times faster by working modulo $p$ and $q$ separately and gluing the results with the \emph{Chinese remainder theorem}. Every public-key operation is a modular power, computed by \emph{square-and-multiply}. Finally, Diffie--Hellman and ElGamal rely on the fact that undoing a modular power, the \emph{discrete logarithm}, is infeasible.

\subsection{Primes and how many there are}
A positive integer is \imp{prime if it has exactly two divisors, 1 and itself}; a \emph{composite} has more than two divisors, and the number 1 (exactly one divisor) is neither. Two integers are \emph{coprime} (relatively prime) when their gcd is 1, so a prime $p$ is coprime to every integer from 1 to $p - 1$, which is why $\mathbb{Z}_p^* = \{1, \dots, p - 1\}$.

\textbf{Infinitely many primes.} Suppose the primes were only $p_1, \dots, p_k$, and let $N = p_1 p_2 \cdots p_k + 1$. Dividing $N$ by any $p_i$ leaves remainder 1, so $N$ either is a new prime or has a prime factor outside the list; either way the list was incomplete. Note that $N$ itself need not be prime: $2 \cdot 3 \cdot 5 \cdot 7 \cdot 11 \cdot 13 + 1 = 30031 = 59 \times 509$, but 59 and 509 are new primes.

\textbf{Counting primes.} The function $\pi(n)$ counts the primes $\le n$: $\pi(1) = 0$, $\pi(2) = 1$, $\pi(3) = 2$, $\pi(10) = 4$, $\pi(20) = 8$, $\pi(50) = 15$, $\pi(100) = 25$. For large $n$ there is no formula, only bounds (Lagrange found the lower one, Gauss the upper):
\[
\frac{n}{\ln n} \;<\; \pi(n) \;<\; \frac{n}{\ln n - 1.08366}.
\]
For $n = 10^6$: $n/\ln n = 72{,}382.4$ and $n/(\ln n - 1.08366) = 78{,}543.2$, so $72{,}383 \le \pi(10^6) \le 78{,}543$; the actual count is $78{,}498$. The practical consequence is that primes thin out only slowly: near $N$ about one number in $\ln N$ is prime. For 1024-bit numbers $\ln 2^{1024} \approx 710$, so if we test random \emph{odd} 1024-bit numbers, about one in 355 is prime, and generating an RSA prime takes a few hundred primality tests. This is why the test itself must be fast.

\textbf{Checking a number by $\sqrt n$.} If $n = ab$ with $1 < a \le b$, then $a^2 \le ab = n$, so $a \le \sqrt n$. Therefore a composite $n$ always has a prime factor $\le \lfloor \sqrt n \rfloor$, and it is enough to try the primes up to $\lfloor\sqrt n\rfloor$.
\begin{itemize}
  \item Is 97 prime? $\lfloor\sqrt{97}\rfloor = 9$, and none of 2, 3, 5, 7 divides 97, so \textbf{97 is prime}.
  \item Is 301 prime? $\lfloor\sqrt{301}\rfloor = 17$, so we would try 2, 3, 5, 7, 11, 13, 17; but already $7 \mid 301$ ($301 = 7 \times 43$), so \textbf{301 is composite}.
\end{itemize}

\textbf{Mersenne and Fermat numbers.} These were early attempts at formulas that produce primes. A \emph{Mersenne number} is $M_p = 2^p - 1$. If the exponent factors as $ab$, then $2^a - 1$ divides $2^{ab} - 1$ (since $x - 1 \mid x^b - 1$ with $x = 2^a$), so $M_p$ can be prime only for prime $p$. The condition is not sufficient: $M_2, M_3, M_5, M_7, M_{13}, M_{17}, M_{19} = 3, 7, 31, 127, 8191, 131071, 524287$ are prime, but $M_{11} = 2047 = 23 \times 89$ and $M_{23} = 8388607 = 47 \times 178481$. A \emph{Fermat number} is $F_n = 2^{2^n} + 1$. Fermat saw that $F_0, \dots, F_4 = 3, 5, 17, 257, 65537$ are prime and guessed that all are, but Euler found $F_5 = 4294967297 = 641 \times 6700417$, and no larger Fermat number is known to be prime. $F_4 = 65537 = 2^{16} + 1$ survives in practice as the standard RSA public exponent, because its binary form has only two 1-bits, which makes encryption with square-and-multiply cheap.

\subsection{Primality testing: deterministic and probabilistic}
The problem is to decide, for a number of several hundred digits, whether it is prime. A \emph{deterministic} algorithm always gives the correct answer; a \emph{probabilistic} algorithm is much faster but may declare a composite number ``prime'' with a small, controllable probability. Its ``composite'' answers are always correct, because they come with a witness.

\textbf{Divisibility test (deterministic).} Try every $r = 2, 3, \dots, \lfloor\sqrt n\rfloor$ and return ``composite'' as soon as $r \mid n$, otherwise ``prime''. It is correct by the $\sqrt n$ argument above, but its cost is exponential in the size of the input: an $n_b$-bit number has $\sqrt n \approx 2^{n_b/2}$ candidate divisors. For a 200-bit $n$ this is $2^{100}$ bit operations; at $2^{30}$ operations per second that takes $2^{70}$ seconds, about $3.7 \times 10^{13}$ years.

\textbf{AKS (deterministic, polynomial).} In 2002 Agrawal, Kayal and Saxena gave the first deterministic test whose time is polynomial in the number of bits, $O((\log_2 n)^{12})$ bit operations, based on the identity $(x - a)^p \equiv x^p - a \pmod p$, which holds exactly when $p$ is prime. It settled a theoretical question, but in practice the probabilistic tests below are far faster.

\textbf{Fermat test (probabilistic).} Fermat's little theorem says that if $n$ is prime then $a^{n-1} \equiv 1 \pmod n$ for every $a$ with $n \nmid a$. Read in reverse: \imp{if $a^{n-1} \not\equiv 1 \pmod n$, then $n$ is certainly composite}; if $a^{n-1} \equiv 1$, then $n$ is only \emph{probably} prime, because some composites also satisfy it. For example $341 = 11 \times 31$ gives $2^{340} \equiv 1 \pmod{341}$ (since $2^{10} = 1024 = 3 \cdot 341 + 1$), so base 2 is fooled, but base 3 exposes it: $3^{340} \equiv 56 \pmod{341}$.

\textbf{Carmichael numbers.} Trying more bases does not always help. A \emph{Carmichael number} is a composite $n$ with $a^{n-1} \equiv 1 \pmod n$ for \emph{every} $a$ coprime to $n$, so the Fermat test fails for all such bases. The smallest is $561 = 3 \times 11 \times 17$ (then 1105, 1729, 2465, 2821). The reason is visible from the factors: $561 - 1 = 560$ is divisible by $3 - 1$, $11 - 1$ and $17 - 1$, so for each prime factor $p$, $a^{560} = (a^{p-1})^{560/(p-1)} \equiv 1 \pmod p$, and a number that is $\equiv 1$ modulo 3, 11 and 17 is $\equiv 1$ modulo 561 (by the CRT below). In general $n$ is Carmichael exactly when it is squarefree and $(p - 1) \mid (n - 1)$ for every prime $p \mid n$ (Korselt's criterion). Only bases sharing a factor with $n$ catch it, e.g.\ $3^{560} \equiv 375 \pmod{561}$, and for large $n$ such bases are too rare to find by chance.

\textbf{Square-root test.} If $n = p$ is prime and $x^2 \equiv 1 \pmod p$, then $p \mid (x - 1)(x + 1)$, and since $p$ is prime it divides one of the two factors, so $x \equiv \pm 1 \pmod p$. Hence \imp{modulo a prime, the only square roots of 1 are $+1$ and $-1$} (written $n - 1$ for $-1$). A composite modulus usually has extra, \emph{nontrivial} square roots of 1, and finding one proves that $n$ is composite.
\begin{center}\small
\begin{tabular}{@{}lll@{}}
\toprule
$n$ & Square roots of 1 mod $n$ & Comment\\
\midrule
7 (prime) & 1, 6 & only $\pm1$\\
17 (prime) & 1, 16 & only $\pm1$\\
8 (composite) & 1, 3, 5, 7 & extra roots 3, 5 $\Rightarrow$ composite\\
15 (composite) & 1, 4, 11, 14 & extra roots 4, 11\\
22 (composite) & 1, 21 & only $\pm1$ although composite\\
561 (composite) & 1, 67, 188, 254, 307, 373, 494, 560 & $2^3 = 8$ roots (3 prime factors)\\
\bottomrule
\end{tabular}
\end{center}
The test is one-way: $n = 22 = 2 \times 11$ has only $\pm 1$, because modulo 2 there is a single square root of 1. For an odd $n$ with $r$ distinct prime factors there are $2^r$ square roots of 1 (a choice of $\pm1$ modulo each prime factor, glued by the CRT). A nontrivial root also \emph{factors} $n$: since $n \mid (x - 1)(x + 1)$ but divides neither factor, $\gcd(x \pm 1, n)$ are proper factors; for $x = 67$, $\gcd(66, 561) = 33$ and $\gcd(68, 561) = 17$.

\subsection{Miller--Rabin test \new}
\likely{``Use the Miller--Rabin test with base 2 to check whether 561 (or 61, 341) is prime'': write $n - 1 = m \cdot 2^k$, then the $T$ table.}
The Fermat test only looks at the final value $a^{n-1}$, and Carmichael numbers make that value 1 for every base. Miller--Rabin also looks at the values \emph{on the way} to $a^{n-1}$ and applies the square-root test to them, which catches Carmichael numbers too.

\textbf{Idea.} Write $n - 1 = m \cdot 2^k$ with $m$ odd ($n$ odd, so $k \ge 1$). Then $a^{n-1}$ is obtained from $a^m$ by squaring $k$ times:
\[
a^m,\; a^{2m},\; a^{4m},\; \dots,\; a^{2^k m} = a^{n-1} \pmod n,
\]
each term being the square of the previous one. If $n$ is prime, the last term is 1 by Fermat. Walking backwards from it, every term whose square is 1 must be $\pm1$, by the square-root test. So for a prime $n$ the sequence either starts with $a^m \equiv 1$ (then it is all 1s), or it contains a $-1$ somewhere before the last position, after which it is all 1s. If the sequence reaches $+1$ \emph{without} a $-1$ just before it, the previous term is a nontrivial square root of 1 and $n$ is composite; if it never reaches $\pm1$, Fermat already fails.

\begin{mustknow}[: Miller--Rabin (base $a$)]
Find $m$ odd and $k$ with $n - 1 = m \times 2^k$. Set $T \leftarrow a^m \bmod n$. If $T = \pm1$, return \textbf{prime}.\\
Repeat at most $k - 1$ times: $T \leftarrow T^2 \bmod n$; if $T = +1$ return \textbf{composite} (the previous $T$ was a nontrivial square root of 1); if $T = -1$ return \textbf{prime}.\\
If the loop ends without a decision, return \textbf{composite}.
\end{mustknow}
The loop stops after $k - 1$ squarings because the $k$-th square is $a^{n-1}$ itself: if the previous $T$ is not $\pm1$, then either its square is 1 (a nontrivial root) or it is not 1 (Fermat fails), so $n$ is composite in both cases and nothing is gained by computing it. ``Prime'' here means ``probably prime'': a composite $n$ passes for at most $1/4$ of the bases (such bases are called \emph{strong liars}), so $t$ independent random bases leave an error probability of at most $4^{-t}$, e.g.\ $4^{-40} \approx 10^{-24}$. The \emph{recommended primality test} (Forouzan) is therefore trial division by small primes first, which cheaply removes most composites, followed by Miller--Rabin with several bases.

\begin{worked}[: Miller--Rabin on 561 and 61 (Forouzan, base 2)]
\textbf{Q.} Use the Miller--Rabin test with base 2 to decide whether (a) 561 and (b) 61 are prime.

\textbf{(a)} $n - 1 = 560 = 35 \times 2^4$, so $m = 35$, $k = 4$, $a = 2$, and at most $k - 1 = 3$ squarings.
First $2^{35} \bmod 561$ by repeated squaring: $2^8 = 256$, $2^{16} = 256^2 \bmod 561 = 460$, $2^{32} = 460^2 \bmod 561 = 103$, so $2^{35} = 2^{32} \cdot 2^3 \equiv 103 \times 8 = 824 \equiv 263$.
\begin{center}\small
\begin{tabular}{@{}llll@{}}
\toprule
Step & Computation & $T$ & Decision\\
\midrule
Init & $2^{35} \bmod 561$ & 263 & not $\pm1$, continue\\
1 & $263^2 = 69169 \bmod 561$ & 166 & not $\pm1$, continue\\
2 & $166^2 = 27556 \bmod 561$ & 67 & not $\pm1$, continue\\
3 & $67^2 = 4489 \bmod 561$ & $\mathbf{1}$ & $+1$ $\Rightarrow$ \textbf{composite}\\
\bottomrule
\end{tabular}
\end{center}
\textbf{561 is composite.} Justification: $67^2 \equiv 1$ with $67 \ne \pm1$, so 67 is a nontrivial square root of 1, impossible modulo a prime. (Indeed $561 = 3 \times 11 \times 17$; and $\gcd(66, 561) = 33$, $\gcd(68, 561) = 17$.) Note that one more squaring gives $2^{560} \equiv 1$, so 561 \emph{passes} the Fermat test; Miller--Rabin catches it only because it inspects the step from 67 to 1.

\medskip
\textbf{(b)} $n - 1 = 60 = 15 \times 2^2$, so $m = 15$, $k = 2$, at most one squaring. $2^{15} = 2^8 \cdot 2^4 \cdot 2^2 \cdot 2 \equiv 12 \times 16 \times 4 \times 2 = 1536 = 25 \cdot 61 + 11 \equiv 11 \pmod{61}$.
\begin{center}\small
\begin{tabular}{@{}llll@{}}
\toprule
Step & Computation & $T$ & Decision\\
\midrule
Init & $2^{15} \bmod 61$ & 11 & not $\pm1$, continue\\
1 & $11^2 = 121 \bmod 61$ & $60 = \mathbf{-1}$ & $-1$ $\Rightarrow$ \textbf{prime}\\
\bottomrule
\end{tabular}
\end{center}
\textbf{61 passes (probably prime)}, and 61 is indeed prime. Check: $2^{60} = (2^{30})^2 \equiv (-1)^2 = 1$, consistent with Fermat.
\end{worked}

\begin{worked}[: Fermat test vs Miller--Rabin on 341 (new)]
\textbf{Q.} Show that 341 passes the Fermat test with base 2, then apply the Miller--Rabin test with base 2. What can you conclude?

\textbf{Fermat test.} $2^{10} = 1024 = 3 \times 341 + 1$, so $2^{10} \equiv 1 \pmod{341}$ and $2^{340} = (2^{10})^{34} \equiv 1$. 341 passes the Fermat test with base 2.

\textbf{Miller--Rabin.} $340 = 85 \times 2^2$, so $m = 85$, $k = 2$, at most one squaring. $2^{85} = (2^{10})^{8} \cdot 2^5 \equiv 1 \cdot 32 = 32$.
\begin{center}\small
\begin{tabular}{@{}llll@{}}
\toprule
Step & Computation & $T$ & Decision\\
\midrule
Init & $2^{85} \bmod 341$ & 32 & not $\pm1$, continue\\
1 & $32^2 = 1024 \bmod 341$ & $\mathbf{1}$ & $+1$ $\Rightarrow$ \textbf{composite}\\
\bottomrule
\end{tabular}
\end{center}
\textbf{341 is composite}: 32 is a nontrivial square root of 1 modulo 341. Base 2 is a Fermat liar for 341 but not a strong liar. The root even gives the factors: $\gcd(31, 341) = 31$ and $\gcd(33, 341) = 11$, so $341 = 11 \times 31$.
\end{worked}

\begin{worked}[: a prime with a long sequence (new)]
\textbf{Q.} Apply the Miller--Rabin test with base 2 to $n = 97$.

$96 = 3 \times 2^5$, so $m = 3$, $k = 5$, at most 4 squarings.
\begin{center}\small
\begin{tabular}{@{}llll@{}}
\toprule
Step & Computation & $T$ & Decision\\
\midrule
Init & $2^3 \bmod 97$ & 8 & continue\\
1 & $8^2 = 64$ & 64 & continue\\
2 & $64^2 = 4096 = 42 \cdot 97 + 22$ & 22 & continue\\
3 & $22^2 = 484 = 5 \cdot 97 - 1$ & $96 = \mathbf{-1}$ & \textbf{prime}\\
\bottomrule
\end{tabular}
\end{center}
\textbf{97 passes} (and is prime, as the $\sqrt{97}$ check showed earlier).
\end{worked}

The remaining slide runs and two instructive edge cases:
\begin{center}\small
\begin{tabular}{@{}llll@{}}
\toprule
$n$, base & $n - 1 = m \cdot 2^k$ & $T$ sequence & Result\\
\midrule
27, 2 & $13 \cdot 2^1$ & 11 (no loop, $k - 1 = 0$) & composite (correct)\\
4033, 2 & $63 \cdot 2^6$ & 3521, \textbf{4032} ($= -1$) & ``prime'' (wrong: $37 \times 109$)\\
4033, 3 & $63 \cdot 2^6$ & 3551, 2443, 3442, 2443, 3442, 2443 & composite (correct)\\
2047, 2 & $1023 \cdot 2^1$ & \textbf{1} & ``prime'' (wrong: $M_{11} = 23 \times 89$)\\
25, 7 & $3 \cdot 2^3$ & 18, \textbf{24} ($= -1$) & ``prime'' (wrong: $5^2$)\\
\bottomrule
\end{tabular}
\end{center}
Bases 2 for 4033 and 2047, and 7 for 25, are strong liars; this is exactly why the test is repeated with several bases (4033 fools base 2 but not base 3).

\subsection{Factorization}
\textbf{Fundamental theorem of arithmetic.} Every integer $n > 1$ can be written as a product of primes, $n = p_1^{e_1} p_2^{e_2} \cdots p_k^{e_k}$, and this representation is unique apart from the order of the factors. Uniqueness lets us read gcd and lcm directly from the exponents. A common divisor can use each prime at most as often as \emph{both} numbers contain it, and a common multiple must use it at least as often as \emph{either} does:
\[
\gcd(a, b) = \prod p_i^{\min(a_i, b_i)}, \qquad \operatorname{lcm}(a, b) = \prod p_i^{\max(a_i, b_i)}, \qquad \gcd(a, b) \times \operatorname{lcm}(a, b) = a \times b,
\]
the last because $\min(a_i, b_i) + \max(a_i, b_i) = a_i + b_i$ for every prime. Example: $2160 = 2^4 \cdot 3^3 \cdot 5$ and $1200 = 2^4 \cdot 3 \cdot 5^2$ give $\gcd = 2^4 \cdot 3 \cdot 5 = 240$ and $\operatorname{lcm} = 2^4 \cdot 3^3 \cdot 5^2 = 10800$; check $240 \times 10800 = 2{,}592{,}000 = 2160 \times 1200$. This method is useless for large numbers, since it needs the factorization, which is exactly the hard part; Euclid's algorithm finds the gcd without it.

\textbf{Why factorization matters.} Finding $\phi(n)$, and with it the RSA private key, is as hard as factoring $n = pq$. Primality testing is easy (Miller--Rabin), but finding the factors of a number already known to be composite is believed to be hard, and RSA lives in the gap between the two problems.

\textbf{Trial division.} Divide $n$ by $a = 2, 3, 4, \dots$ while $a \le \sqrt n$; whenever $a \mid n$, output $a$, replace $n$ by $n/a$, and try the same $a$ again (to catch repeated factors). If $n > 1$ when the loop ends, the remaining $n$ is a prime factor. Examples: $1233 = 3^2 \times 137$, and $1523357784 = 2^3 \times 3^2 \times 13 \times 37 \times 43987$. Like the divisibility test, it costs up to $\sqrt n = 2^{n_b/2}$ divisions, so it only works when $n$ has small factors.

\textbf{Fermat's factorization method.} Any odd $n = ab$ (with $a \ge b$) is a difference of two squares: with $x = (a + b)/2$ and $y = (a - b)/2$ we get $x^2 - y^2 = ab = n$. Conversely, if we find $x$ with $w = x^2 - n$ a perfect square $y^2$, then $n = (x + y)(x - y)$. The method starts at $x = \lceil\sqrt n\rceil$ and increases $x$ by 1 until $x^2 - n$ is a square. Since the search ends at $x = (a + b)/2$, it is fast exactly when the two factors are close to $\sqrt n$.
\begin{center}\small
\begin{tabular}{@{}cccl@{}}
\toprule
$x$ & $w = x^2 - 5959$ & Square? & \\
\midrule
78 & 125 & no & ($\lceil\sqrt{5959}\rceil = 78$)\\
79 & 282 & no & \\
80 & 441 & yes, $21^2$ & $5959 = (80 + 21)(80 - 21) = \mathbf{101 \times 59}$\\
\bottomrule
\end{tabular}
\end{center}
For $8051$ it succeeds immediately ($90^2 - 8051 = 49 = 7^2$, so $8051 = 97 \times 83$), while $4033 = 37 \times 109$ needs 10 steps ($x = 64, \dots, 73$; $73^2 - 4033 = 1296 = 36^2$). The security lesson is that RSA primes $p$ and $q$ must not be close to each other.

\textbf{Other methods} (named on the slides). \emph{Pollard $p - 1$} finds a prime factor $p$ quickly when $p - 1$ has only small prime factors, by computing $\gcd(a^{B!} - 1, n)$. \emph{Pollard rho} iterates $x \leftarrow x^2 + c \bmod n$ and finds a factor $p$ in about $\sqrt p$ steps from a repetition modulo $p$, detected with a gcd. The \emph{quadratic sieve} and the \emph{number field sieve} build many congruences $x^2 \equiv y^2 \pmod n$ (a generalisation of Fermat's idea) and are sub-exponential; the number field sieve is the fastest known method for large RSA moduli, but it is still far too slow for 2048-bit $n$. (A \emph{sieve of Eratosthenes}, by contrast, is the school method for \emph{listing} all primes up to a bound by crossing out multiples.)

\subsection{Chinese remainder theorem \new}
\likely{Forouzan style ``Solve $x \equiv 2 \pmod 3$, $x \equiv 3 \pmod 5$, $x \equiv 2 \pmod 7$'' with the full $M$, $M_i$, $M_i^{-1}$ table and the inverses by extended Euclid; or a word problem / RSA speed-up.}
\textbf{Problem.} We are given several congruences in one unknown with different moduli,
\[
x \equiv a_1 \pmod{m_1},\quad x \equiv a_2 \pmod{m_2},\quad \dots,\quad x \equiv a_k \pmod{m_k},
\]
and want all $x$ that satisfy them simultaneously. \imp{The Chinese remainder theorem (CRT) states that if the moduli are pairwise coprime, the system has exactly one solution modulo $M = m_1 m_2 \cdots m_k$}, and it gives a formula for it. For example $x \equiv 2 \pmod 3$, $x \equiv 3 \pmod 5$, $x \equiv 2 \pmod 7$ has the solution $x = 23$ (and $23 + 105t$).

\textbf{Idea.} Build $x$ as a sum of $k$ terms in which the $i$-th term is $\equiv a_i$ modulo $m_i$ and $\equiv 0$ modulo every other modulus. Then modulo $m_i$ all terms but the $i$-th vanish and $x \equiv a_i$, for each $i$ at once. A number divisible by every modulus except $m_i$ is $M_i = M / m_i$; multiplying it by its inverse modulo $m_i$ turns it into 1 modulo $m_i$ without disturbing the zeros, and multiplying by $a_i$ gives the required value.

\begin{mustknow}[: CRT in four steps]
1.\ $M = m_1 \times m_2 \times \cdots \times m_k$ (the common modulus).\quad
2.\ $M_i = M / m_i$ for each $i$.\\
3.\ $M_i^{-1}$ = the inverse of $M_i$ \textbf{modulo $m_i$} (reduce $M_i \bmod m_i$ first, then extended Euclid).\\
4.\ $x = \left(a_1 M_1 M_1^{-1} + a_2 M_2 M_2^{-1} + \cdots + a_k M_k M_k^{-1}\right) \bmod M$.
\end{mustknow}

\textbf{Why it works.} For $j \ne i$, $m_j$ divides $M_i$, so the term $a_i M_i M_i^{-1}$ is $\equiv 0 \pmod{m_j}$. Modulo $m_i$ the term equals $a_i \cdot (M_i M_i^{-1}) \equiv a_i \cdot 1$. Therefore $x \equiv a_i \pmod{m_i}$ for every $i$. Step 3 needs $\gcd(M_i, m_i) = 1$, and this is exactly where pairwise coprimality is used: $M_i$ is a product of moduli each coprime to $m_i$.

\textbf{Why the solution is unique modulo $M$.} If $x$ and $x'$ both solve the system, every $m_i$ divides $x - x'$; since the $m_i$ are pairwise coprime, their product $M$ divides $x - x'$. So all solutions are $x + tM$, and there is exactly one in $\{0, \dots, M - 1\}$. Equivalently, the map $x \mapsto (x \bmod m_1, \dots, x \bmod m_k)$ is a one-to-one correspondence between $\mathbb{Z}_M$ and $\mathbb{Z}_{m_1} \times \cdots \times \mathbb{Z}_{m_k}$, which is what the applications below exploit.

\textbf{Edge case: moduli not coprime.} Then the theorem does not apply, and the system may have no solution or several. For instance $x \equiv 1 \pmod 4$ and $x \equiv 2 \pmod 6$ has no solution, because the first forces $x$ odd and the second forces $x$ even.

\begin{worked}[: the standard CRT question (Forouzan, slide example)]
\textbf{Q.} Solve $x \equiv 2 \pmod 3$, $x \equiv 3 \pmod 5$, $x \equiv 2 \pmod 7$ using the Chinese remainder theorem.

\textbf{Step 1.} The moduli 3, 5, 7 are pairwise coprime, so the CRT applies. $M = 3 \times 5 \times 7 = 105$.

\textbf{Steps 2--3.}
\begin{center}\small
\begin{tabular}{@{}ccccccc@{}}
\toprule
$i$ & $a_i$ & $m_i$ & $M_i = M/m_i$ & $M_i \bmod m_i$ & $M_i^{-1} \bmod m_i$ & $a_i M_i M_i^{-1}$\\
\midrule
1 & 2 & 3 & 35 & 2 & 2 & $2 \cdot 35 \cdot 2 = 140$\\
2 & 3 & 5 & 21 & 1 & 1 & $3 \cdot 21 \cdot 1 = 63$\\
3 & 2 & 7 & 15 & 1 & 1 & $2 \cdot 15 \cdot 1 = 30$\\
\bottomrule
\end{tabular}
\end{center}
Inverses (extended Euclid with $r_1 = m_i$, $r_2 = M_i \bmod m_i$, $t_1 = 0$, $t_2 = 1$). For $35^{-1} \bmod 3 = 2^{-1} \bmod 3$:
\begin{center}\small
\begin{tabular}{@{}rrrrrrr@{}}
$q$ & $r_1$ & $r_2$ & $r$ & $t_1$ & $t_2$ & $t$\\ \midrule
1 & 3 & 2 & 1 & 0 & 1 & $-1$\\
2 & 2 & 1 & 0 & 1 & $-1$ & 3\\
 & \textbf{1} & 0 & & $\mathbf{-1}$ & 3 & \\
\end{tabular}
\end{center}
$r_1 = 1$, so the inverse exists and equals $t_1 = -1 \equiv \mathbf{2} \pmod 3$ (check $2 \cdot 2 = 4 \equiv 1$). Since $21 \equiv 1 \pmod 5$ and $15 \equiv 1 \pmod 7$, the other two inverses are 1.

\textbf{Step 4.} $x = (140 + 63 + 30) \bmod 105 = 233 \bmod 105 = \mathbf{23}$.

\textbf{Check.} $23 = 7 \cdot 3 + 2$, $23 = 4 \cdot 5 + 3$, $23 = 3 \cdot 7 + 2$ \checkmark. All solutions: $x = 23 + 105t$.
\end{worked}

\begin{worked}[: CRT with non-trivial inverses (new)]
\textbf{Q.} Find the smallest positive $x$ with $x \equiv 3 \pmod 7$, $x \equiv 5 \pmod{11}$, $x \equiv 7 \pmod{13}$.

\textbf{Step 1.} 7, 11, 13 are distinct primes, hence pairwise coprime. $M = 7 \times 11 \times 13 = 1001$.

\textbf{Steps 2--3.}
\begin{center}\small
\begin{tabular}{@{}ccccccc@{}}
\toprule
$i$ & $a_i$ & $m_i$ & $M_i$ & $M_i \bmod m_i$ & $M_i^{-1} \bmod m_i$ & $a_i M_i M_i^{-1}$\\
\midrule
1 & 3 & 7 & 143 & 3 & 5 & $3 \cdot 143 \cdot 5 = 2145$\\
2 & 5 & 11 & 91 & 3 & 4 & $5 \cdot 91 \cdot 4 = 1820$\\
3 & 7 & 13 & 77 & 12 & 12 & $7 \cdot 77 \cdot 12 = 6468$\\
\bottomrule
\end{tabular}
\end{center}
Inverses by extended Euclid (columns $q, r_1, r_2, r, t_1, t_2, t$):
\begin{center}\small
\begin{tabular}{@{}rrrrrrr@{\qquad}rrrrrrr@{}}
\multicolumn{7}{@{}l}{$3^{-1} \bmod 7$} & \multicolumn{7}{l@{}}{$3^{-1} \bmod 11$}\\
\midrule
$q$ & $r_1$ & $r_2$ & $r$ & $t_1$ & $t_2$ & $t$ & $q$ & $r_1$ & $r_2$ & $r$ & $t_1$ & $t_2$ & $t$\\
2 & 7 & 3 & 1 & 0 & 1 & $-2$ & 3 & 11 & 3 & 2 & 0 & 1 & $-3$\\
3 & 3 & 1 & 0 & 1 & $-2$ & 7 & 1 & 3 & 2 & 1 & 1 & $-3$ & 4\\
 & \textbf{1} & 0 & & $\mathbf{-2}$ & 7 & & 2 & 2 & 1 & 0 & $-3$ & 4 & $-11$\\
 & & & & & & & & \textbf{1} & 0 & & \textbf{4} & $-11$ & \\
\end{tabular}
\end{center}
So $143^{-1} \equiv 3^{-1} \equiv -2 \equiv \mathbf{5} \pmod 7$ (check $3 \cdot 5 = 15 \equiv 1$) and $91^{-1} \equiv 3^{-1} \equiv \mathbf{4} \pmod{11}$ (check $3 \cdot 4 = 12 \equiv 1$). For the third, $77 \equiv 12 \equiv -1 \pmod{13}$, and $(-1)^{-1} = -1$, so $77^{-1} \equiv \mathbf{12} \pmod{13}$.

\textbf{Step 4.} $x = (2145 + 1820 + 6468) \bmod 1001 = 10433 \bmod 1001 = \mathbf{423}$.

\textbf{Check.} $423 = 60 \cdot 7 + 3$, $\;423 = 38 \cdot 11 + 5$, $\;423 = 32 \cdot 13 + 7$ \checkmark.
\end{worked}

\begin{trap}[: inverses modulo the wrong number]
$M_i^{-1}$ is the inverse modulo $m_i$ (the small modulus), not modulo $M$. Reduce $M_i$ modulo $m_i$ before running extended Euclid, but multiply the \emph{unreduced} $M_i$ in step 4; using $M_i \bmod m_i$ in the final sum breaks the zeros modulo the other moduli.
\end{trap}

\textbf{Congruences with a coefficient, $c\,x \equiv b \pmod m$.} The CRT formula needs every congruence in the form $x \equiv a_i \pmod{m_i}$, so each one is first solved for $x$. If $\gcd(c, m) = 1$, multiply both sides by $c^{-1} \bmod m$; the modulus does not change. If $d = \gcd(c, m) > 1$, the congruence has a solution only when $d \mid b$, and then dividing \emph{everything} by $d$ gives $\frac{c}{d}x \equiv \frac{b}{d} \pmod{\frac{m}{d}}$: the modulus shrinks to $m/d$, because $m \mid c(x - x')$ only forces $\frac{m}{d} \mid (x - x')$. Dividing $c$ and $b$ by $d$ but keeping the old modulus $m$ silently throws away $d - 1$ of the $d$ solutions modulo $m$.

\begin{worked}[: CRT after reducing the coefficients (video example 2, corrected)]
\textbf{Q.} Solve $4x \equiv 5 \pmod 9$ and $2x \equiv 6 \pmod{20}$.

\textbf{First congruence.} $\gcd(4, 9) = 1$ and $4^{-1} \bmod 9 = 7$ (since $4 \cdot 7 = 28 = 3 \cdot 9 + 1$). Multiplying by 7: $x \equiv 35 \equiv 8 \pmod 9$.

\textbf{Second congruence.} $\gcd(2, 20) = 2$ and $2 \mid 6$, so it is solvable; dividing by 2 gives $x \equiv 3 \pmod{10}$, which modulo 20 means $x \equiv 3$ \emph{or} $13 \pmod{20}$.

\textbf{CRT on $x \equiv 8 \pmod 9$, $x \equiv 3 \pmod{10}$.} $\gcd(9, 10) = 1$, $M = 90$.
\begin{center}\small
\begin{tabular}{@{}ccccccc@{}}
\toprule
$i$ & $a_i$ & $m_i$ & $M_i$ & $M_i \bmod m_i$ & $M_i^{-1} \bmod m_i$ & $a_i M_i M_i^{-1}$\\
\midrule
1 & 8 & 9 & 10 & 1 & 1 & $8 \cdot 10 \cdot 1 = 80$\\
2 & 3 & 10 & 9 & 9 & 9 & $3 \cdot 9 \cdot 9 = 243$\\
\bottomrule
\end{tabular}
\end{center}
$x = (80 + 243) \bmod 90 = 323 \bmod 90 = \mathbf{53}$, so the complete solution is $\mathbf{x \equiv 53 \pmod{90}}$, i.e.\ $x \in \{53, 143\}$ modulo 180.

\textbf{Check.} $4 \cdot 53 = 212 = 23 \cdot 9 + 5$ and $2 \cdot 53 = 106 = 5 \cdot 20 + 6$ \checkmark; $4 \cdot 143 = 572 = 63 \cdot 9 + 5$ and $2 \cdot 143 = 286 = 14 \cdot 20 + 6$ \checkmark.

\emph{The video} kept modulus 20 after dividing ($x \equiv 3 \pmod{20}$) and used $M = 180$, $M_1^{-1} = 20^{-1} \bmod 9 = 5$, $M_2^{-1} = 9^{-1} \bmod 20 = 9$, getting $x = 1043 \bmod 180 = 143$. That value is correct but it is only one of the two solutions modulo 180; the branch $x \equiv 13 \pmod{20}$ gives the other, 53.
\end{worked}

\textbf{More slide examples} (all checked):
\begin{center}\small
\begin{tabular}{@{}lllll@{}}
\toprule
System & $M$ & $M_i$ & $M_i^{-1}$ & $x$\\
\midrule
$x \equiv 3\ (7),\ 3\ (13),\ 0\ (12)$ & 1092 & 156, 84, 91 & 4, 11, 7 & \textbf{276}\\
$S \equiv 1\ (11),\ 12\ (13),\ 2\ (17)$ & 2431 & 221, 187, 143 & 1, 8, 5 & \textbf{155}\\
$z \equiv 61\ (99),\ 65\ (98),\ 69\ (97)$ & 941094 & 9506, 9603, 9702 & 50, 97, 49 & \textbf{457}\\
\bottomrule
\end{tabular}
\end{center}

\textbf{Applications.}
\begin{itemize}
  \item \emph{Word problems.} ``Find an integer that leaves remainder 3 when divided by 7 and by 13 but is divisible by 12'' is the first row above: $x = 276$ ($276 = 12 \times 23$, $276 = 39 \cdot 7 + 3 = 21 \cdot 13 + 3$).
  \item \emph{Secret sharing.} The secret $S$ is handed out as residues modulo 11, 13, 17, 19: shares 1, 12, 2, 3. Any 3 shares determine $S$ uniquely modulo the product of their moduli, which exceeds $S$, so CRT recovers it; every one of the four 3-subsets gives $S = 155$ (the second row uses $\{1, 12, 2\}$).
  \item \emph{Arithmetic on big numbers with small registers.} To compute $z = 123 + 334$ on a machine that holds only numbers below 100, represent $123 \equiv 24, 25, 26$ and $334 \equiv 37, 40, 43$ modulo $99, 98, 97$; add componentwise to get $z \equiv 61, 65, 69$; CRT gives $z = 457$, correct because $457 < M = 941094$ and the solution is unique modulo $M$.
  \item \emph{RSA decryption speed-up.} Since $n = pq$ with $p, q$ coprime, the plaintext $P = C^d \bmod n$ is determined by $P \bmod p$ and $P \bmod q$. These are computed with exponents reduced by Fermat, $d_p = d \bmod (p - 1)$ and $d_q = d \bmod (q - 1)$, on numbers half as long, and are then combined by CRT. A modular multiplication on half-length numbers costs about a quarter as much and the exponents are half as long, so the two half-size powers together are roughly four times faster than one full-size power.
\end{itemize}

\begin{worked}[: RSA decryption with CRT (new)]
\textbf{Q.} In RSA with $p = 7$, $q = 11$, $n = 77$, $e = 13$, $d = 37$, decrypt $C = 41$ using the CRT.

\textbf{Reduced exponents.} $d_p = 37 \bmod 6 = 1$, \;$d_q = 37 \bmod 10 = 7$.

\textbf{Half-size powers.} $P_p = 41^{1} \bmod 7 = 6$.\\
$P_q = 41^{7} \bmod 11 = 8^7 \bmod 11$: $8^2 = 64 \equiv 9$, $8^4 \equiv 81 \equiv 4$, $8^7 = 8^4 \cdot 8^2 \cdot 8 \equiv 4 \cdot 9 \cdot 8 = 288 \equiv 2$.

\textbf{CRT.} Solve $P \equiv 6 \pmod 7$, $P \equiv 2 \pmod{11}$ with $M = 77$:\\
$M_1 = 11 \equiv 4 \pmod 7$, $4^{-1} \bmod 7 = 2$; \;$M_2 = 7$, $7^{-1} \bmod 11 = 8$ ($7 \cdot 8 = 56 = 5 \cdot 11 + 1$).\\
$P = (6 \cdot 11 \cdot 2 + 2 \cdot 7 \cdot 8) \bmod 77 = (132 + 112) \bmod 77 = 244 \bmod 77 = \mathbf{13}$.

\textbf{Check.} Directly, $41^{37} \bmod 77 = 13$, and re-encrypting, $13^{13} \bmod 77 = 41$ \checkmark.
\end{worked}

\subsection{Fast exponentiation: square-and-multiply \hot}
\likely{Needed inside every RSA / ElGamal / Diffie--Hellman numerical: show the bit table for $a^x \bmod n$.}
\textbf{Problem.} Public-key cryptography computes $y = a^x \bmod n$ with $x$ and $n$ of 1024 bits or more. Multiplying by $a$ again and again needs $x - 1$ multiplications, about $2^{1024}$, which is impossible. (The size of intermediate results is the smaller problem: reducing modulo $n$ after every multiplication keeps them below $n^2$.)

\textbf{Idea.} Write the exponent in binary, $x = x_{n_b - 1} 2^{n_b - 1} + \cdots + x_1 2^1 + x_0 2^0$ with each $x_i \in \{0, 1\}$. Then
\[
a^x = \prod_{i \,:\, x_i = 1} a^{2^i}, \qquad\text{e.g.}\quad a^{22} = a^{16} \cdot a^{4} \cdot a^{2} \quad (22 = 10110_2).
\]
The powers $a, a^2, a^4, a^8, \dots$ are each the square of the previous one, so they cost one squaring each, and we multiply into the result only those whose bit is 1.

\textbf{Mechanism (Forouzan, right-to-left).} Start with $y = 1$. Scan the bits from $x_0$ (least significant) upwards: if $x_i = 1$, set $y \leftarrow a \times y \bmod n$; then square, $a \leftarrow a^2 \bmod n$ (not needed after the last bit). Return $y$. An $n_b$-bit exponent needs at most $n_b - 1$ squarings and $n_b$ multiplications, so the cost is $O(n_b)$ modular multiplications, linear in the number of bits instead of exponential. A 1024-bit exponent needs at most about 2048 multiplications. The public exponent $e = 65537 = 10000000000000001_2$ needs 16 squarings and a single non-trivial multiplication, which is why it is the usual choice.

\begin{center}\small
\textbf{$17^{22} \bmod 21$ (slide example), $22 = 10110_2$}\\[2pt]
\begin{tabular}{@{}ccll@{}}
\toprule
$i$ & $x_i$ & Multiplication (init.\ $y = 1$) & Squaring (init.\ $a = 17$)\\
\midrule
0 & 0 & -- & $a = 17^2 \bmod 21 = 16$\\
1 & 1 & $y = 1 \times 16 \bmod 21 = 16$ & $a = 16^2 \bmod 21 = 4$\\
2 & 1 & $y = 16 \times 4 \bmod 21 = 1$ & $a = 4^2 \bmod 21 = 16$\\
3 & 0 & -- & $a = 16^2 \bmod 21 = 4$\\
4 & 1 & $y = 1 \times 4 \bmod 21 = \mathbf{4}$ & (not needed)\\
\bottomrule
\end{tabular}
\end{center}
So $17^{22} \bmod 21 = 4$, using 4 squarings and 3 multiplications instead of 21 multiplications. Section~\ref{sec:modexp} computes the same kind of power left-to-right (scanning from the most significant bit, square then multiply by $a$); both give the same answer, and either is acceptable in the exam as long as every intermediate value is shown.

\begin{worked}[: RSA encryption and decryption by square-and-multiply (Forouzan numbers)]
\textbf{Q.} In RSA, $n = 77$, $e = 13$, $d = 37$. Encrypt $M = 5$ and verify by decrypting, using the square-and-multiply method.

\textbf{Encryption} $C = 5^{13} \bmod 77$, $13 = 1101_2$ ($x_0 = 1, x_1 = 0, x_2 = 1, x_3 = 1$):
\begin{center}\small
\begin{tabular}{@{}ccll@{}}
\toprule
$i$ & $x_i$ & Multiplication ($y = 1$) & Squaring ($a = 5$)\\
\midrule
0 & 1 & $y = 1 \times 5 = 5$ & $a = 5^2 = 25$\\
1 & 0 & -- & $a = 25^2 = 625 \bmod 77 = 9$\\
2 & 1 & $y = 5 \times 9 = 45$ & $a = 9^2 = 81 \bmod 77 = 4$\\
3 & 1 & $y = 45 \times 4 = 180 \bmod 77 = \mathbf{26}$ & (not needed)\\
\bottomrule
\end{tabular}
\end{center}
\textbf{Decryption} $M = 26^{37} \bmod 77$, $37 = 100101_2$:
\begin{center}\small
\begin{tabular}{@{}ccll@{}}
\toprule
$i$ & $x_i$ & Multiplication ($y = 1$) & Squaring ($a = 26$)\\
\midrule
0 & 1 & $y = 26$ & $a = 26^2 = 676 \bmod 77 = 60$\\
1 & 0 & -- & $a = 60^2 = 3600 \bmod 77 = 58$\\
2 & 1 & $y = 26 \times 58 = 1508 \bmod 77 = 45$ & $a = 58^2 = 3364 \bmod 77 = 53$\\
3 & 0 & -- & $a = 53^2 = 2809 \bmod 77 = 37$\\
4 & 0 & -- & $a = 37^2 = 1369 \bmod 77 = 60$\\
5 & 1 & $y = 45 \times 60 = 2700 \bmod 77 = \mathbf{5}$ & (not needed)\\
\bottomrule
\end{tabular}
\end{center}
\textbf{Answer.} $C = 26$, and decryption returns $M = 5$ \checkmark.
\end{worked}

\subsection{Discrete logarithms \new}
\likely{``Find $x$ such that $3^x \equiv 7 \pmod{17}$'' via a table of powers, using the log rules; or theory: why the discrete log is hard and which systems rely on it.}
\textbf{Problem.} Square-and-multiply computes $y = a^x \bmod n$ quickly. The reverse question is: given $a$, $y$ and $n$, find $x$. Over the real numbers this is $\log_a y$; modulo $n$ it is the \emph{discrete logarithm}, and no comparably fast method is known for it.

\textbf{Primitive roots (reminder).} The order of $a$ in $\mathbb{Z}_n^*$ is the smallest $i \ge 1$ with $a^i \equiv 1$, and $a$ is a \emph{primitive root} when its order equals $\phi(n)$; details and examples are in \S\ref{sec:fermat}. For a prime $p$, $\mathbb{Z}_p^*$ always has primitive roots, and if $g$ is one then $g^1, g^2, \dots, g^{p-1}$ are all different and run through every element $1, \dots, p - 1$ exactly once. Hence every $y \in \mathbb{Z}_p^*$ equals $g^x$ for exactly one $x \in \{1, \dots, p - 1\}$; this $x$ is the \emph{discrete logarithm} of $y$ to the base $g$, written $x = L_g\, y$ (the modulus is understood). In $\mathbb{Z}_7^*$ the primitive roots are 3 and 5, giving two logarithm tables:
\begin{center}\small
\begin{tabular}{@{}c|cccccc@{}}
$y$ & 1&2&3&4&5&6\\ \midrule
$x = L_3\, y$ (mod 7) & 6&2&1&4&5&3\\
$x = L_5\, y$ (mod 7) & 6&4&5&2&1&3\\
\end{tabular}
\end{center}
With the table, $4 \equiv 3^x \pmod 7$ gives $x = L_3 4 = 4$ ($3^4 = 81 \equiv 4$), and $6 \equiv 5^x \pmod 7$ gives $x = L_5 6 = 3$ ($5^3 = 125 \equiv 6$).

\textbf{Rules.} Since $g^{p-1} \equiv 1$, exponents of $g$ only matter modulo $p - 1 = \phi(p)$, so discrete logs obey the usual logarithm laws with arithmetic modulo $\phi(n)$:
\[
L_g 1 \equiv 0, \qquad L_g(x \cdot y) \equiv L_g x + L_g y, \qquad L_g(x^k) \equiv k \cdot L_g x \pmod{\phi(n)}.
\]
The second rule follows from $g^{L x} \cdot g^{L y} = g^{L x + L y}$. (The table writes $L_g 1 = 6$ rather than 0, because it lists exponents from 1 to $p - 1$; $6 \equiv 0 \pmod 6$.) Example in $\mathbb{Z}_7^*$: $L_3(4 \cdot 5) = L_3 6 = 3$, and $L_3 4 + L_3 5 = 4 + 5 = 9 \equiv 3 \pmod 6$ \checkmark.

\textbf{Edge case: base not a primitive root.} Then the powers of $a$ do not cover the whole group, so $a^x \equiv y$ may have no solution or several. In $\mathbb{Z}_7^*$, 2 has order 3 ($2^x$ cycles through 2, 4, 1), so $2^x \equiv 3 \pmod 7$ has no solution, while $2^x \equiv 4$ has $x = 2$ and $x = 5$.

\textbf{Why it is hard.} The only general method on the slides is \emph{exhaustive search}: for $x = 1, 2, \dots, n - 1$, test whether $y = a^x \bmod n$ and return the first $x$ that works. That is up to $p - 1$ exponentiations, exponential in the number of bits, while computing $a^x$ needs only $O(\log p)$ multiplications. Better algorithms exist (baby-step giant-step in about $\sqrt p$ steps; index calculus, sub-exponential), but for a well-chosen 2048-bit prime they are still infeasible, roughly as hard as factoring a 2048-bit RSA modulus. Modular exponentiation is therefore a \emph{one-way function}, and this asymmetry is the basis of \textbf{Diffie--Hellman} (both parties publish $g^a$ and $g^b$ and share $g^{ab}$; an eavesdropper would need $a = L_g(g^a)$) and of \textbf{ElGamal} (the private key $d$ hides behind the public $e_2 = e_1^d \bmod p$). For example, with $p = 17$, $g = 3$, $a = 4$, $b = 7$: the public values are $3^4 \equiv 13$ and $3^7 \equiv 11$, and the shared key is $11^4 \equiv 13^7 \equiv 4 \pmod{17}$.

\begin{worked}[: discrete logarithms modulo 17 (new)]
\textbf{Q.} (a) Show that 3 is a primitive root of 17 and tabulate $L_3\, y$ for $\mathbb{Z}_{17}^*$. (b) Solve $3^x \equiv 7 \pmod{17}$. (c) Using only $L_3 2$ and $L_3 7$ and the rules, find $L_3 14$. (d) Solve $3^x \equiv 8 \pmod{17}$ from $L_3 2$.

\textbf{(a)} Successive powers $3^x \bmod 17$, each $3 \times$ the previous:
\begin{center}\small
\setlength{\tabcolsep}{3.5pt}
\begin{tabular}{@{}c|cccccccccccccccc@{}}
$x$ & 1&2&3&4&5&6&7&8&9&10&11&12&13&14&15&16\\ \midrule
$3^x \bmod 17$ & 3&9&10&13&5&15&11&16&14&8&7&4&12&2&6&1\\
\end{tabular}
\end{center}
The first power equal to 1 is $3^{16}$, so $\operatorname{ord}(3) = 16 = \phi(17)$ and 3 is a primitive root (all 16 values differ). Reading the table backwards:
\begin{center}\small
\setlength{\tabcolsep}{3.5pt}
\begin{tabular}{@{}c|cccccccccccccccc@{}}
$y$ & 1&2&3&4&5&6&7&8&9&10&11&12&13&14&15&16\\ \midrule
$L_3\, y$ & 16&14&1&12&5&15&11&10&2&3&7&13&4&9&6&8\\
\end{tabular}
\end{center}
\textbf{(b)} $x = L_3 7 = \mathbf{11}$. Check: $3^{11} \bmod 17 = 7$ (from the power table).

\textbf{(c)} $L_3 14 = L_3(2 \cdot 7) \equiv L_3 2 + L_3 7 = 14 + 11 = 25 \equiv \mathbf{9} \pmod{16}$. Check: $3^9 \bmod 17 = 14$ \checkmark.

\textbf{(d)} $L_3 8 = L_3(2^3) \equiv 3 \cdot L_3 2 = 42 \equiv \mathbf{10} \pmod{16}$. Check: $3^{10} \bmod 17 = 8$ \checkmark.
\end{worked}

\begin{fixbox}
\begin{itemize}
  \item Slide (Fermat test): ``$561 = 33 \times 17$'' is true but is not the prime factorisation; $561 = 3 \times 11 \times 17$, the smallest Carmichael number.
  \item Slide (AKS): the example computes $(\log_2 200)^{12} \approx 3.9 \times 10^{10}$ bit operations for a 200-bit $n$ (the printed 39,547,615,483 is not even that value exactly; it is $\approx 39{,}787{,}815{,}242$). For a 200-bit number $\log_2 n = 200$, so the bound $(\log_2 n)^{12}$ is $200^{12} \approx 4.1 \times 10^{27}$: still polynomial, but not ``40 seconds''.
  \item Slide (recommended test, 4033): the list of trial divisors ``2, 3, 5, 7, 11, 17, 23'' skips 13 and 19 (neither divides 4033). A full divisibility test up to $\lfloor\sqrt{4033}\rfloor = 63$ would find 37 and settle the question; the recommended test uses trial division by \emph{small} primes only.
  \item Slide (Fermat test): ``If $x$ is a prime, then $a^{n-1} \equiv 1$'' should read ``if $n$ is a prime''.
  \item Slide $\phi(1) = 0$ follows Forouzan's ``integers smaller than $n$''; the standard convention is $\phi(1) = 1$. It never matters in RSA. (Totient chapter.)
  \item Euler's second version $a^{k\phi(n)+1} \equiv a \pmod n$ holds for squarefree $n$ (such as $n = pq$ in RSA), not for every $n$; it fails for $n = 4$, $a = 2$. (See \S\ref{sec:fermat}.)
\end{itemize}
\end{fixbox}
